\chapter{Roadmap}

A evolução do projeto foi organizada em um roadmap que permite visualizar o estado atual e os próximos passos de forma objetiva. A tabela a seguir resume as áreas de evolução e seu status de implementação no momento desta documentação.

\begin{table}[H]
    \centering
    \caption{Roadmap arquitetural e técnico do PlantDiseaseAI}
    \label{tab:roadmap}
    \begin{tabular}{p{4cm}p{4cm}p{5cm}}
        \toprule
        \textbf{Área} & \textbf{Status} & \textbf{Observações} \\
        \midrule
        Infraestrutura & Parcialmente concluída & API, Storage local, Repository em memória e mensageria configurada \\
        IA & Em preparação & Pipeline placeholder, ModelLoader e contratos prontos \\
        Banco & Não implementado & Persistência temporária em memória \\
        Docker & Não iniciado & Aguardando consolidação da infraestrutura \\
        Kubernetes & Não iniciado & Dependente da estratégia de containerização \\
        Observabilidade & Parcialmente planejada & Logs estruturados podem ser expandidos \\
        Autenticação & Não implementada & Fora do escopo atual \\
        \bottomrule
    \end{tabular}
\end{table}

O roadmap indica que a base estrutural já suporta a introdução do modelo treinado sem alterações profundas. As próximas fases deverão focar na substituição dos placeholders por componentes produtivos, na preparação da execução em container e na definição de observabilidade e segurança conforme a necessidade de implantação.
